\documentclass{standalone}
\usepackage{circuitikz}
\usetikzlibrary{calc}
\definecolor{circuitnoteink}{HTML}{000000}
\definecolor{circuitnotemark}{HTML}{FE00FE}
\begin{document}
\begin{circuitikz}[american, line width=0.8pt]
\ctikzset{bipoles/length=1.2cm}
\coordinate (c2) at (2,-4);
\coordinate (c4) at (6,-4);
\coordinate (c5) at (8,-4);
\coordinate (c7) at (12,-4);
\coordinate (c8) at (14,-4);
\coordinate (c10) at (18,-4);
\coordinate (b3) at (4,-2);
\coordinate (b6) at (10,-2);
\coordinate (b9) at (16,-2);
\draw (c2) to[potentiometer, n=part-P1, l_=$P_{1}$, a^=$10\,\mathrm{k}\Omega$] (c4); % line 3
\draw (c5) to[thyristor, n=part-T1, l_=$T_{1}$] (c7); % line 4
\draw (c8) to[triac, n=part-T2, l_=$T_{2}$] (c10); % line 5
\draw (part-P1.wiper) -- (b3); % line 7
\draw (part-T1.gate) |- (b6); % line 8
\draw (part-T2.gate) |- (b9); % line 9
\draw[circuitnoteink] (1,-0.2) -- (19,-0.2); % line 11
\draw[circuitnoteink] (1,-8.2) -- (19,-8.2); % line 12
\draw[circuitnoteink] (1,-0.2) -- (1,-8.2); % line 13
\draw[circuitnoteink] (7,-0.2) -- (7,-8.2); % line 14
\draw[circuitnoteink] (13,-0.2) -- (13,-8.2); % line 15
\draw[circuitnoteink] (19,-0.2) -- (19,-8.2); % line 16
\path (1.841,-1.393) rectangle (6.159,-0.207); % line 17
\node[anchor=west, inner sep=0, circuitnotemark, font=\large] at (4,-0.8) {X}; % line 17
\path (0.571,-7.993) rectangle (7.429,-6.807); % line 18
\node[anchor=west, inner sep=0, circuitnotemark, font=\large] at (4,-7.4) {X}; % line 18
\path (8.508,-1.393) rectangle (11.492,-0.207); % line 19
\node[anchor=west, inner sep=0, circuitnotemark, font=\large] at (10,-0.8) {X}; % line 19
\path (7.587,-7.993) rectangle (12.413,-6.807); % line 20
\node[anchor=west, inner sep=0, circuitnotemark, font=\large] at (10,-7.4) {X}; % line 20
\path (14.286,-1.393) rectangle (17.715,-0.207); % line 21
\node[anchor=west, inner sep=0, circuitnotemark, font=\large] at (16,-0.8) {X}; % line 21
\path (13.968,-7.993) rectangle (18.032,-6.807); % line 22
\node[anchor=west, inner sep=0, circuitnotemark, font=\large] at (16,-7.4) {X}; % line 22
\node[anchor=south west, inner sep=2pt, circuitnotemark, font=\LARGE\bfseries] (circuittitle) at (current bounding box.north west) {X};
\path (circuittitle.south west) rectangle ++(12.314,0.853);
\end{circuitikz}
\end{document}
